\chapter{System Setup and Commissioning}
\label{ch:setup}

Before the impedance controller can be executed on the Franka Emika Panda,
the complete hardware and software environment must be correctly configured.
This chapter covers the setup of the workstation, network connection,
Franka Desk and FCI access, real-time Ubuntu kernel, libfranka installation,
and C++ build environment, as specified in Objective~2 of the thesis
registration.

\section{Hardware Architecture and Network Connection}
\label{sec:hw_network}

The experimental platform consists of the Panda arm, its Franka Control
Unit (FCU), a workstation PC, and an emergency stop device.
The workstation is connected \textbf{directly} to the FCI Ethernet port
of the FCU with a Cat\,6a cable---no intermediate switch is permitted,
as any additional hop introduces latency sufficient to trigger an FCI
communication timeout.

Static IP addresses are assigned:

\begin{center}
\begin{tabular}{@{}lll@{}}
\toprule
Device & Interface & Address \\
\midrule
FCU (FCI port) & \texttt{eth0} & \texttt{172.16.0.2} \\
Workstation    & \texttt{eth1} & \texttt{172.16.0.1/24} \\
\bottomrule
\end{tabular}
\end{center}

Connectivity is verified with \texttt{ping 172.16.0.2}; the round-trip
time must be consistently below $1\,\mathrm{ms}$.

\section{Franka Desk and FCI Access}
\label{sec:franka_desk}

Franka Desk is the web-based configuration interface, accessible
at \url{https://172.16.0.2} from any browser on the workstation.
Before real-time torque commands can be sent, the following steps are
required:

\begin{enumerate}
  \item Unlock the robot joints via the Franka Desk \textit{Joints} page.
  \item Activate FCI mode in \textit{Settings} $\to$ \textit{Network}.
  \item Confirm that no error is shown on the status indicator.
  \item Place the emergency stop within reach of the operator.
\end{enumerate}

All experiments are conducted with FCI active and the emergency stop
accessible at all times.

\section{Real-Time Ubuntu Kernel}
\label{sec:rt_kernel}

The FCI requires that the control callback completes within 1\,ms.
A standard Linux kernel cannot guarantee this because high-priority threads
can be preempted by interrupt handlers.
The PREEMPT\_RT patch makes all kernel threads and interrupt handlers fully
preemptible, reducing worst-case scheduling latency to below $50\,\mu s$.

Installation steps:

\begin{enumerate}
  \item Install Ubuntu 22.04 LTS on the workstation.
  \item Download and apply the PREEMPT\_RT patch for the corresponding
        kernel version.
  \item Compile and install the patched kernel:
        \texttt{make -j\$(nproc) deb-pkg}.
  \item Grant real-time scheduling priority to the control user:
        add \texttt{@realtime soft rtprio 99} to
        \texttt{/etc/security/limits.conf}.
  \item Set CPU governor to performance:
        \texttt{cpufreq-set -g performance}.
\end{enumerate}

The maximum latency of the patched system is verified with
\texttt{cyclictest}, which must show a worst-case latency below
$200\,\mu s$ under load before any FCI experiments are conducted.

\section{libfranka Installation}
\label{sec:libfranka}

\texttt{libfranka} is built from source to match the firmware version of
the robot controller:

\begin{lstlisting}[style=cppstyle, language=bash,
  caption={Build and install libfranka from source},
  label={lst:libfranka}]
sudo apt install build-essential cmake git libpoco-dev libeigen3-dev
git clone --recursive https://github.com/frankaemika/libfranka
cd libfranka && git checkout 0.9.2
mkdir build && cd build
cmake -DCMAKE_BUILD_TYPE=Release -DBUILD_TESTS=OFF ..
cmake --build . -j$(nproc)
sudo make install
\end{lstlisting}

The installation is verified by running the
\texttt{communication\_test} example program provided with libfranka,
which must report zero packet losses over a 60-second run.

\section{C++ Build Environment}
\label{sec:cpp_env}

All controller code is built with CMake in Release mode (\texttt{-O3}),
which provides approximately three times faster matrix operations compared
to Debug mode.
The following libraries are required:

\begin{itemize}
  \item \textbf{Eigen\,3}: header-only linear algebra library for all
        matrix and vector operations.
  \item \textbf{libfranka}: robot communication and model interface.
  \item \textbf{Poco}: network and threading utilities (dependency of
        libfranka).
\end{itemize}

A minimal \texttt{CMakeLists.txt}:

\begin{lstlisting}[style=cppstyle, language=bash,
  caption={Minimal CMakeLists.txt for a libfranka controller},
  label={lst:cmake}]
cmake_minimum_required(VERSION 3.10)
project(impedance_controller)
set(CMAKE_CXX_STANDARD 14)
set(CMAKE_BUILD_TYPE Release)
find_package(Franka REQUIRED)
find_package(Eigen3 REQUIRED)
add_executable(impedance_ctrl src/impedance_controller.cpp)
target_link_libraries(impedance_ctrl Franka::Franka Eigen3::Eigen)
\end{lstlisting}

%% =============================================================
%%  CHAPTER 4: SOFTWARE IMPLEMENTATION
%% =============================================================
